\clearpage
\VerdierPhysicalPageAnchor{0191}%
\begingroup\small\setlength{\parskip}{0pt}\setlength{\emergencystretch}{1.5em}
\setlength{\abovedisplayskip}{4pt plus 2pt minus 2pt}%
\setlength{\belowdisplayskip}{4pt plus 2pt minus 2pt}%

\VerdierSemanticAnchor{ch3d-2-2-2}
\noindent\textbf{補題 2.2.2.}\quad
\textit{$0\longrightarrow X\longrightarrow Y\longrightarrow Z\longrightarrow0$
を $\mathcal A$ における完全列とする。}

\noindent\textit{a) $X$ と $Y$ が $M^{\wedge}$ の元ならば、エピ射
$u\colon L\longrightarrow Z$ が存在し、$L\in M$ である。}

\noindent\textit{b) $X$ と $Z$ が $M^{\wedge}$ の元ならば、エピ射
$u\colon L\longrightarrow Y$ が存在し、$L\in M$ である。}

\noindent\textit{c) $Y$ と $Z$ が $M^{\wedge}$ の元ならば、エピ射
$u\colon L\longrightarrow X$ が存在し、$L\in M$ である。}

\VerdierSemanticAnchor{ch3d-2-2-3}
\medskip\noindent\textbf{補題 2.2.3.}\quad
\textit{$0\longrightarrow X\longrightarrow Y\longrightarrow Z\longrightarrow0$
を $\mathcal A$ における完全列とする。}

\noindent\textit{a) $Y$ と $Z$ が $M^{\wedge}$ の元ならば、任意のエピ射
$u\colon L\longrightarrow X$（$L\in M$）に対して、$u$ の核は
$M^{\wedge}$ の二つの元の間のあるエピ射の核である。}

\noindent\textit{b) $X$ と $Y$ が $M^{\wedge}$ の元ならば、任意のエピ射
$u\colon L\longrightarrow Z$（$L\in M$）に対して、$u$ の核は
$M^{\wedge}$ の二つの元の間のあるモノ射の余核である。}

\noindent\textit{c) $X$ と $Z$ が $M^{\wedge}$ の元ならば、任意のエピ射
$u\colon L\longrightarrow Y$（$L\in M$）に対して、$u$ の核は
$M^{\wedge}$ の二つの元の間のあるエピ射の核である。}

\VerdierSemanticAnchor{ch3d-2-2-2-3-proof}
二つの補題の証明は読者に委ねる。補題 (2.2.2) と (2.2.3) の性質
(a), (c) では性質 $(E_2)$ を用いないことに注意せよ。

\VerdierSemanticAnchor{ch3d-2-2-4}
\medskip\noindent\textbf{命題 2.2.4.}\quad
\textit{a) 集合 $M^{\wedge}$ は性質 $(E_6)$ と $(E_2)$ をもつ。集合
$M$ は $M^{\wedge}$ に含まれる。$M^{\wedge}$ の各対象は $M$ のある対象の商である。}

\noindent\textit{b) 対象の集合 $M'$ が $M$ を含み、$M'$ の各対象が $M$ の
ある対象の商であり、かつ性質 $(E_3)$ をもつならば、その集合は
$M^{\wedge}$ に含まれる。}

\noindent\textit{c) 特に、
$(M^{\wedge})^{\wedge}=M^{\wedge}$ である。}

\VerdierSemanticAnchor{ch3d-2-2-4-proof}
性質 a) は直前の二つの補題から従う。性質 b), c) は明らかである。

\VerdierSemanticAnchor{ch3d-2-2-5}
\medskip\noindent\textbf{2.2.5.}\quad
$M\subset M'$ を $\mathcal A$ の対象の二つの集合とし、$M$ は性質
$(E_4)$ と $(E_5)$ を、$M'$ は性質 $(E_4)$ をもち、さらに $M'$ の各対象は
$M$ のある対象の商であるとする。対象 $X$（$M'$ に属する）の $M$-次元が
$n$ 以下であるとは、次の完全列が存在することをいい、
$M\text{-}\dim(X)\leq n$ と記す：

\VerdierSemanticAnchor{ch3d-2-2-5-exact-sequence}
\[
0\longrightarrow L_n\longrightarrow L_{n-1}\longrightarrow\cdots
\longrightarrow L_0\longrightarrow X\longrightarrow0,
\qquad L_i\in M,\quad0\leq i\leq n.
\]
\endgroup
